\section{Lecture 18: Normalization Methods}

\subsection{Motivation}

In the previous lectures, we examined two key aspects of training deep networks:

\begin{itemize}
    \item gradient instability (vanishing and exploding gradients),
    \item architectural choices such as activation functions and depth.
\end{itemize}

\medskip

Even with careful initialization and suitable activations, training deep networks can remain unstable.

\medskip

\textbf{Key question.}  
How can we stabilize the distribution of activations during training?

\medskip

\textbf{Guiding idea.}  
Normalization methods control the scale and distribution of intermediate representations.

\subsection{Batch Normalization}

Batch normalization is one of the most influential techniques in modern deep learning.

\medskip

Consider activations in a layer:
\[
z = W x + b.
\]

Batch normalization normalizes these activations across a mini-batch.

\medskip

For a batch $\{z_1,\dots,z_m\}$:

\[
\mu_B = \frac{1}{m}\sum_{i=1}^m z_i
\]

\[
\sigma_B^2 = \frac{1}{m}\sum_{i=1}^m (z_i-\mu_B)^2
\]

Normalized activations:

\[
\hat{z}_i = \frac{z_i - \mu_B}{\sqrt{\sigma_B^2 + \epsilon}}
\]

To preserve representational flexibility:

\[
y_i = \gamma \hat{z}_i + \beta
\]

where $\gamma$ and $\beta$ are learnable parameters.

\medskip

\textbf{Interpretation}

\begin{itemize}
\item normalize mean and variance
\item then allow the network to re-scale if needed
\end{itemize}

\medskip

\textbf{Benefits}

\begin{itemize}
\item stabilizes activation distributions
\item accelerates training
\item reduces sensitivity to initialization
\end{itemize}

\subsection{Layer and Instance Normalization}

Batch normalization relies on mini-batch statistics, which may not always be appropriate.

\medskip

Alternative normalization schemes include:

\textbf{Layer normalization}

Normalize across features of a single example:

\[
\mu = \frac{1}{d}\sum_{j=1}^{d} z_j
\]

\[
\sigma^2 = \frac{1}{d}\sum_{j=1}^{d}(z_j-\mu)^2
\]

\medskip

\textbf{Instance normalization}

Normalize each example independently across spatial dimensions (commonly used in vision tasks).

\medskip

\textbf{Comparison}

\begin{itemize}
\item Batch normalization: uses batch statistics
\item Layer normalization: uses feature statistics
\item Instance normalization: uses per-sample spatial statistics
\end{itemize}

\medskip

\textbf{Practical observation}

Layer normalization is commonly used in transformer architectures.

\subsection{Stabilization of Training}

Normalization improves training dynamics in several ways.

\medskip

\textbf{Gradient stabilization}

\begin{itemize}
\item keeps activations within a reasonable range
\item prevents gradients from exploding or vanishing
\end{itemize}

\medskip

\textbf{Improved conditioning}

\begin{itemize}
\item reduces sensitivity to parameter scaling
\item improves optimization geometry
\end{itemize}

\medskip

\textbf{Regularization effect}

\begin{itemize}
\item stochasticity from batch statistics can act as implicit regularization
\end{itemize}

\medskip

\textbf{Insight}

Normalization reshapes the optimization landscape.

\subsection{Theoretical and Practical Insights}

Although batch normalization was originally motivated heuristically, several explanations have been proposed.

\medskip

\textbf{Variance control}

Normalization prevents signal magnitudes from growing or shrinking across layers.

\medskip

\textbf{Optimization smoothing}

Empirical studies suggest normalization makes the loss landscape smoother and easier to optimize.

\medskip

\textbf{Gradient flow}

Normalized activations help maintain stable gradient magnitudes during backpropagation.

\medskip

\textbf{Architectural interaction}

Normalization often works together with other design choices:

\begin{itemize}
\item residual connections
\item ReLU activations
\item careful initialization
\end{itemize}

\subsection{A Unifying Perspective}

Normalization methods address a central challenge of deep learning:

\begin{itemize}
\item controlling the propagation of signals through deep networks.
\end{itemize}

They interact with:

\begin{itemize}
\item activation functions,
\item optimization methods,
\item architectural design.
\end{itemize}

\medskip

\textbf{Core idea}

Stable training requires controlling both forward activations and backward gradients.

\subsection{Outlook}

In this lecture we introduced:

\begin{itemize}
\item batch normalization,
\item layer and instance normalization,
\item the role of normalization in training stability.
\end{itemize}

\medskip

These techniques allow much deeper networks to be trained successfully.

\medskip

In the next lecture, we study regularization methods in deep learning, including dropout and data augmentation, which improve generalization in large neural networks.

\medskip

\textbf{Guiding principle}

Deep learning success depends on carefully balancing expressivity, optimization, and regularization.